\section{Введение}
\label{sec:Chapter0} \index{Chapter0}

В современных сетях возникает необходимость в гибкой обработке пакетов. Необходимо быстро внедрять новые протоколы и сервисы без замены оборудования. Решением этого вопроса является использование программируемых сетевых устройств с архитектурой RMT с языком описания пакетной обработки P4. Такие устройства реализуют конвейерную обработку, состоящую из последовательных match‑action стадий.

При компиляции P4-программы для архитектуры RMT нужно распределять логические таблицы по стадиям конвейера. От этого решения зависят число задействованных стадий и, следовательно, задержка обработки пакета. Для успешного размещения следует учитывать ограничения архитектуры аппаратуры и порядок применения таблиц, заданный графом зависимостей таблиц (Table Dependency Graph). На практике стремятся минимизировать число занятых стадий, поскольку оно непосредственно определяет задержку обработки пакета.

В данной работе рассматриваются жадные алгоритмы решения задачи размещения, анализируется их эффективность и быстрота выполнения, и предлагается модификация с поддержкой горизонтального разрезания таблиц. В главе 2 приводится описание предметной области, в главе 3 формулируется цель и основные задачи работы , в главе 4 описываются реализованные алгоритмы, в главе 5 приводятся результаты экспериментального исследования, в главе 6 формулируются выводы.



